%%%%%%%%%%%%%%%%%%%%%%%%%%%%%%%%%%%%%%%%%%%%%%%%%%%%%%%%%%%%%%%%%%%%%%
% Lớp 9 | Chương I. Phương trình và hệ hai phương trình bậc nhất hai ẩn | Giải hệ hai phương trình bậc nhất hai ẩn
% Dạng 9D1?2-1: Giải hệ bằng phương pháp thế
%%%%%%%%%%%%%%%%%%%%%%%%%%%%%%%%%%%%%%%%%%%%%%%%%%%%%%%%%%%%%%%%%%%%%%


% nguồn: 2026-11-15_09_de_tham_khao_quan_3_de_1_va_de_2 (nhập 15/11/2026)
\begin{bt}%[9D1H2-1]%[ma:68475a0a]
% nguon: Đề tham khảo Tuyển sinh 10 TP.HCM 2025-2026 – Quận 3 – Đề 1 – Bài 4
Một tổ may gồm cả công nhân nam và nữ được giao may $360$ chiếc áo cổ vũ. Để hoàn thành nhiệm vụ, mỗi công nhân nam may $6$ chiếc áo, mỗi công nhân nữ may $8$ chiếc áo. Gọi $x$ và $y$ lần lượt là số công nhân nam và số công nhân nữ ($x, y \in \mathbb{N}^*$).
\begin{enumerate}[a)]
	\item Lập công thức biểu diễn $y$ theo $x$.
	\item Tổ may có bao nhiêu công nhân? Biết rằng số công nhân nữ gấp $1{,}5$ lần số công nhân nam.
\end{enumerate}
\loigiai{
\begin{enumerate}[a)]
	\item Vì $x$ công nhân nam may được $6x$ chiếc áo và $y$ công nhân nữ may được $8y$ chiếc áo, tổng số áo may được là $360$ chiếc nên ta có phương trình:
	$6x + 8y = 360 \Leftrightarrow 8y = 360 - 6x \Leftrightarrow y = \dfrac{360 - 6x}{8} = 45 - \dfrac{3}{4}x$.
	\item Theo giả thiết, số công nhân nữ gấp $1{,}5$ lần số công nhân nam, tức là:
	$y = 1{,}5x = \dfrac{3}{2}x$.\\
	Thay vào phương trình $6x + 8y = 360$, ta được:
	$6x + 8\left(\dfrac{3}{2}x\right) = 360 \Leftrightarrow 6x + 12x = 360 \Leftrightarrow 18x = 360 \Leftrightarrow x = 20\text{ (công nhân)}$.\\
	Suy ra số công nhân nữ là: $y = 1{,}5 \cdot 20 = 30\text{ (công nhân)}$.\\
	Vậy tổ may có tất cả $20 + 30 = 50\text{ công nhân}$.
\end{enumerate}
}
\end{bt}
